\documentclass[11pt,a4paper]{article}
\usepackage[margin=23mm]{geometry}
\usepackage{amsmath,amssymb}
\usepackage[T1]{fontenc}
\usepackage[hidelinks]{hyperref}
\setlength{\parindent}{0pt}
\setlength{\parskip}{6pt}
\newcommand{\R}{\mathbb R}
\title{RM10: Recovery/acceleration, information, and useful ramp gradients}
\author{Methodology review --- Option A only}
\date{5 October 2026}
\begin{document}
\maketitle
\textbf{Conclusion.} Physical acceleration/recovery has support. A positive task
with genuinely hidden, action-relevant information is not admitted in this bounded
review. Conditional aggregate sufficiency and missing rich-family evidence are
different findings. Option C was not opened.

\section*{1. Task and complete population accounting}
This is an audit template, not an adopted traffic update. Let $s=0,\ldots,S$
index a common time grid with step $\Delta t$ (seconds). Let $\theta\in\R^d$
parameterize a bounded common open-loop ramp rate $u_s(\theta)$ (vehicle-equivalents
per second). The state $x_s$ contains road inventories, ramp/external queues and
whatever recovery variables a justified model requires. Public parameters are
$p$ and public arrivals/boundary inputs are $b_s$. A lawful discretization would give
\[
 x_{s+1}=F_s(x_s,u_s(\theta),b_s;p).
\]
Permission $u_s$, release at a meter, and actual entry at a merge are distinct:
queues, travel and receiving can prevent permission being realized. This template
specifies no new release or merge rule.

Let $P_s=P(x_s)$ be \emph{all} arrived but not completed traffic mass, including
external waiting and both road and ramp traffic. Let $a_s$ be external demand
arrival rate and $e_s$ be final scenario exit rate. Internal transfers cancel:
\begin{equation}
 P_{s+1}=P_s+\Delta t(a_s-e_s),\qquad
 J(\theta)=\Delta t\sum_{s=0}^{S-1}P_s(\theta).
 \label{eq:tst}
\end{equation}
$J$ has vehicle-equivalent seconds as units; $J/M$ is mean system time for a
common fixed total demand mass $M$. The grid convention is left-endpoint time
accounting. A common $S$ must clear \emph{all} comparison policies and small
perturbations; truncation with traffic remaining is not full TST.

Iterating the balance gives
\[
 P_s=P_0+\Delta t\sum_{j=0}^{s-1}(a_j-e_j).
\]
For fixed $P_0,a_j$ and fixed $S$, therefore,
\begin{equation}
 \nabla_\theta J
 =-\Delta t^2\sum_{j=0}^{S-2}(S-1-j)\nabla_\theta e_j.
 \label{eq:exit}
\end{equation}
Earlier exits reduce more subsequent waiting intervals. Merely moving a queue
from a mainline lane to the ramp leaves this objective unchanged when final exit
history is unchanged. Mean steady capacity is not the exit-time history.

\newpage
\section*{2. What hidden information must mean}
Let $O$ be a declared observation map with spatial and temporal support. A candidate
history through decision epoch $t$ consists of
\[
 H_t=(O(x_0),\ldots,O(x_t),u_{0:t-1},b_{0:t-1},p).
\]
It includes admitted inventories, completed face counts/flows, speed and ramp
observations. Public future demand may also be common; future observations are
not available before selecting an action.

If $N_i$ is cell mass and actual speed observations are a section mass mean, the
idealized definition is
\[
 N_i=\sum_k m_{ik},\qquad
 U_i=\frac{\sum_k m_{ik}v_{ik}}{N_i},\quad N_i>0.
\]
$m_{ik}$ has vehicle-equivalent units and $v_{ik},U_i$ have speed units; $k$
indexes traffic groups within cell $i$. Empty-cell speed is unavailable.
This illustrates support, not approval of a new sensor contract. Face/time speed
means and equilibrium FD proxies differ from $U_i$. Do not use the closed RM9
FD-speed proxy as an actual non-equilibrium speed law.

Define the consistency set using source-lawful initial states $\mathcal D_0$:
\[
 \mathcal X(H_t)=
 \left\{x_t:\begin{array}{l}
 x_0\in\mathcal D_0,\quad x_{s+1}=F_s(x_s,u_s,b_s;p)\ (0\leq s<t),\\
 O(x_s)=H_t[s]\ (0\leq s\leq t)
 \end{array}\right\}.
\]
Different snapshots with equal total mass do not
establish different future-relevant states in this set. Match public traffic
characteristics and forcing; changing a random future realization is not a
hidden initial-state certificate.

\textbf{Conditional sufficiency.} In a well-defined deterministic model, a known
complete $x_0$ and public inputs determine every $x_s$ by induction. Likewise,
if history identifies a complete state at an earlier epoch, later congestion
does not erase that knowledge. This is conditional on completeness and uniqueness;
an existence result for a PDE alone is not such a uniqueness proof.

For a homogeneous aggregate model with complete state $(N_i,U_i)$ and observed
queues, exact observation observes that complete state. This does not establish
physical adequacy or identifiability of unobserved lane speeds, within-cell waves,
leaders or distributions.

\textbf{A sufficient action certificate.} For two source-lawful states
$x_t^a,x_t^b\in\mathcal X(H_t)$ and two common feasible rate schedules
$\theta_1,\theta_2$, opposite TST rankings demonstrate action relevance:
\[
 J_a(\theta_1)<J_a(\theta_2),\qquad
 J_b(\theta_1)>J_b(\theta_2).
\]
This distinguishes these two actions; it does not prove that all possible optimal
action sets are disjoint. Opposite useful local descent directions could provide
another certificate. Neither is established here.

\newpage
\section*{3. Why recovery alone does not prove necessary retained state}
An inspected density-based aggregate competitor is Srivastava's queue-discharge
demand law (2016 dissertation, Chapter 3.3.1, equation 3.3):
\[
 D_{\rm rec}(k)=\min\{v_fk,\ q_j+c_*(k_j-k)\},\qquad
 c_*=\frac{q_c-q_j}{k_j-k_c}.
\]
$k,k_j,k_c$ are density, jam density and critical density; $v_f$ is free-flow
speed; $q_c$ is capacity; $q_j$ is jam demand; $c_*$ has speed units.
This source changes demand while retaining ordinary supply and conservation.
It is an aggregate recovery precedent for queue discharge, not a validated
ramp-merge law imported by this review.

Another competitor retains a hysteresis mode $\sigma_s$. With public thresholds
$k_-<k_+$ and a known initial mode, a representative recurrence is
\[
 \sigma_{s+1}=
 \begin{cases}
  1,& k_{s+1}\geq k_+,\\
  0,& k_{s+1}\leq k_-,\\
  \sigma_s,& k_-<k_{s+1}<k_+ .
 \end{cases}
\]
This is the known-history logic of Cao/Nilsson/Coogan (2021), Section II.B,
with indexing restated. Complete density history reconstructs this mode.
An unknown initial mode is a different, unresolved domain.

For richer recovery, $N_i$ and $U_i$ constrain only the zeroth and first moments
of a speed distribution $\nu_i$:
\[
 N_i=\int d\nu_i(v),\qquad N_iU_i=\int v\,d\nu_i(v).
\]
These two constraints generally do not specify $\nu_i$. This algebra does
\emph{not} prove that distinct distributions are source-lawful, persist
under the same full history, or change useful ramp-rate actions.

The primary studies support hybrid/local recovery and mean-capacity competitors.
In particular, S12A005's similar mean capacities do not imply exact transient
closure. Its differing transient flow traces are not identical observed
histories. Laurent-Brouty/Costeseque/Goatin's macroscopic model retains local
moving constraints: the label macro cannot be used to dismiss retained state.

\textbf{Completed screening conclusion.} Closed density/rate laws, known
history modes, and completely observed homogeneous speed-state models have
conditional explanations for sufficiency. Rich distribution/leader models
retain physical potential, but lack a joined lawful initialization, equal-history
and useful-action certificate for this task. Negative admission concerns this
review envelope and evidence, not every bounded-acceleration model.

\newpage
\section*{4. General gradient chain and its limits}
The intended path is
\begin{align*}
 \theta&\ \longrightarrow\ {\rm permission/release}
 \longrightarrow\ {\rm actual\ entry/recovery}\\
 &\ \longrightarrow\ {\rm receiving/conservation}
 \longrightarrow\ {\rm final\ exits/TST}.
\end{align*}
Consider a chosen source-lawful discretization on a locally differentiable
branch. Fold the public rate parameterization into
$\widetilde F_s(x,\theta)=F_s(x,u_s(\theta),b_s;p)$. Define
\[
 A_s=\partial_x\widetilde F_s,\quad
 B_s=\partial_\theta\widetilde F_s,\quad
 S_s=\partial_\theta x_s.
\]
For a state of dimension $n$, $A_s\in\R^{n\times n}$ and
$B_s,S_s\in\R^{n\times d}$. If $x_0$ is fixed independently of $\theta$,
\begin{align*}
 S_0&=0,\\
 S_{s+1}&=A_sS_s+B_s,\\
 \nabla_\theta J&=\Delta t\sum_{s=0}^{S-1}
                         S_s^\top\nabla_x P(x_s).
\end{align*}
These are chain-rule identities, not a numerical correctness result.
They agree with the exit-balance expression when both use the same update and
complete population. Normalized gradients divide by the common $M$.

The corresponding adjoint, with no terminal cost and fixed $S$, is
\[
 \lambda_S=0,\qquad
 \lambda_s=\Delta t\nabla_xP(x_s)+A_s^\top\lambda_{s+1},\qquad
 \nabla_\theta J=\sum_{s=0}^{S-1}B_s^\top\lambda_{s+1}.
\]
$\lambda_s\in\R^n$ is the cost sensitivity to $x_s$. Explicit control costs or
control-dependent initialization add their own terms. Neither is introduced here.

Minimum sending/receiving laws, queue caps, active constraints, shocks and
leader interactions can be nonsmooth. Smooth-branch sensitivities do not settle
event boundaries. A hybrid reference may require event-time terms; holding an
event fixed can yield a different derivative. No unapproved smoothing or private
stochastic resampling resolves this issue. An existence theorem is not a
differentiability or useful-gradient theorem. Mean-capacity parameter sensitivities
are different from the ramp-rate TST derivative.

Three evidence levels remain distinct: a justified derivative for a bound task;
numerical/reference agreement; and useful common perturbations in a richer
independent reference. This review gives the conditional general path, not the
latter two. The unresolved history/action contract already prevents positive
core-task admission.

\textbf{Provenance.} Source versions, selected sections, URLs and access limits are
in \path{extraction/rm10/primary_source_audit.md}. The task/consistency definitions,
conditional sufficiency statements and TST chain-rule derivations are reviewer
constructions. They impose no new traffic equation on the repository.
\end{document}
